\section{Proceso de indexación incremental (Modo)}

\subsection{Flujo de procesamiento de documentos nuevos}

LightRAG admite la incorporación dinámica de documentos nuevos; su proceso de Modo (Upsert = Update + Insert) es el siguiente:

\begin{enumerate}
    \item \textbf{Extracción de entidades y relaciones}: con base en un modelo de lenguaje, se extraen entidades (Entity) y relaciones (Edge) del documento nuevo
    \item \textbf{Profiling}: se genera una descripción textual detallada para cada entidad y relación
    \item \textbf{Embedding}: se construye un índice vectorial para las descripciones de entidades, las descripciones de relaciones y el texto de los Chunks
    \item \textbf{Fusión de subgrafos}: el subgrafo del documento nuevo se fusiona con el grafo de conocimiento global
\end{enumerate}

\subsection{Alineación de entidades y consistencia de nombres}

\begin{warningbox}{La consistencia en el nombramiento de entidades es clave}
Una misma entidad puede mencionarse de formas distintas en diferentes Chunks. LightRAG lo especifica explícitamente en el System Prompt de Entity Extraction:
\begin{itemize}
    \item Los nombres de entidades \textbf{no distinguen mayúsculas de minúsculas}
    \item Cada palabra con significado debe escribirse con mayúscula inicial (por ejemplo, «Knowledge Graph» y no «knowledge graph»)
    \item Debe garantizarse un nombramiento consistente (la misma entidad usa siempre el mismo nombre)
\end{itemize}
La alineación de entidades depende principalmente de la consistencia del modelo de lenguaje al generar el texto, y no de un emparejamiento posterior.
\end{warningbox}

\subsection{Estrategia de fusión de descripciones: MapReduce}

Cuando una misma entidad o relación aparece en varios Chunks, cada Chunk aporta una \textbf{descripción parcial} de esa entidad o relación. La estrategia de fusión sigue el patrón MapReduce:

\begin{enumerate}
    \item \textbf{Recolección}: se reúnen todas las descripciones extraídas de los distintos Chunks
    \item \textbf{Evaluación de longitud}:
    \begin{itemize}
        \item Descripción corta $\rightarrow$ se concatena directamente con un separador
        \item Descripción larga $\rightarrow$ se comprime semánticamente (Summarization) con base en un modelo de lenguaje
    \end{itemize}
    \item \textbf{Actualización sincronizada}: \texttt{src\_id} también se concatena, para registrar el origen de la descripción
\end{enumerate}

\begin{importantbox}{Prompt de compresión semántica}
En \texttt{prompts.py} de LightRAG se define el Prompt \texttt{summarize\_entity\_descriptions}, dedicado a comprimir descripciones de entidades o relaciones demasiado extensas en un resumen más conciso, conservando la información esencial.
\end{importantbox}

\subsection{Ejemplo de fusión}

Tomemos dos Chunks como ejemplo:
\begin{itemize}
    \item \textbf{Chunk 1}: se extraen la entidad Alex (descripción A), Bob (descripción A) y la relación Alex→Bob = «colegas»
    \item \textbf{Chunk 2}: se extraen la entidad Alex (descripción B) y la relación Alex→Bob = «amigos»
\end{itemize}

Después de la fusión:
\begin{itemize}
    \item Descripción de Alex = descripción A \texttt{<sep>} descripción B (superposición de descripciones parciales)
    \item Descripción de la relación Alex→Bob = «colegas» \texttt{<sep>} «amigos»
    \item Si la longitud total de la descripción supera el umbral, se activa la compresión semántica
\end{itemize}

\subsection{Resumen de la sección}

La indexación incremental de LightRAG logra la incorporación fluida de documentos nuevos mediante tres pasos: alineación de entidades, fusión de descripciones (MapReduce) y compresión semántica. Todo el proceso depende en gran medida de la capacidad del modelo de lenguaje.
